\chapter{Stand van zaken}\label{ch:stand-van-zaken}In dit hoofdstuk wordt de theoretische fundering van dit onderzoek uiteengezet. De huidige staat van computer vision wordt gekenmerkt door een schaalvergroting die stilaan ecologische, economische en praktische grenzen bereikt. Om deze problematiek adequaat te adresseren, wordt in de volgende secties dieper ingegaan op de wetmatigheden van modelgroei, de beperkingen van bestaande optimalisatiemethoden, en de theoretische werking van onzekerheidskwantificatie als primair filtermechanisme voor datareductie. Dit theoretische kader vormt de noodzakelijke basis voor de methodologie van dit onderzoek, waarin de 'Smart Selection'-methode wordt gedefinieerd.\section{De grenzen van de groei: Scaling Laws}\label{sec:scaling-laws}Binnen deep learning heerst al meer dan een decennium het paradigma dat het vergroten van de modelarchitectuur in combinatie met exponentieel grotere datasets onvermijdelijk leidt tot betere prestaties. Deze \textit{brute-force} benadering heeft geresulteerd in indrukwekkende doorbraken, maar botst inmiddels op fundamentele limieten.\textcite{kaplan2020scaling} hebben dit fenomeen voor het eerst formeel gekwantificeerd in hun zogenaamde \textit{scaling laws}. Zij stelden vast dat de test-foutmarge (loss) van een neuraal netwerk een voorspelbare daling vertoont naarmate de beschikbare rekenkracht ($C$), de datasetgrootte ($D$) en het aantal parameters ($N$) toenemen. Deze relatie is echter niet lineair, maar volgt een machtswet (power-law):$$L(D) \propto D^{-\alpha}$$\textit{Verduidelijking van de formule:} Een korte verduidelijking van deze wiskundige notatie: $L(D)$ staat voor de 'Loss' of de foutmarge van het model, afhankelijk van de hoeveelheid data $D$. Het symbool $\propto$ betekent 'is evenredig met'. De exponent $-\alpha$ (met een cruciaal minteken) geeft aan dat naarmate de hoeveelheid data $D$ stijgt, de foutmarge $L$ daalt. Echter, omdat dit een machtswet is, betekent dit in de praktijk een keiharde wet van verminderde meeropbrengst: om de foutmarge van een model met een vaste, lineaire stap te verbeteren, moet de hoeveelheid data exponentieel groeien.Recentere studies hebben deze wetmatigheden verder verfijnd. \textcite{hoffmann2022training} toonden met de 'Chinchilla'-modellen aan dat veel grote modellen in werkelijkheid ondergetraind zijn ten opzichte van hun grootte. Zij stelden dat voor \textit{compute-optimal} training het aantal modelparameters en het aantal trainingsdata in gelijke verhouding moeten meegroeien. Als een model 10 keer groter wordt gemaakt, moet de dataset ook 10 keer groter worden. Deze ontdekking versterkt de urgentie van datareductie: de onverzadigbare honger naar hoogwaardige trainingsdata is de primaire flessenhals geworden voor de verdere ontwikkeling van kunstmatige intelligentie.\section{Ecologische en economische impact}\label{sec:impact}De exponentiële behoefte aan data en rekenkracht vormt niet enkel een technologische uitdaging, maar vertaalt zich direct in een acuut duurzaamheidsprobleem. De koolstofvoetafdruk van het trainen van state-of-the-art neurale netwerken is inmiddels vergelijkbaar met die van de traditionele zware industrie.In een fundamentele paper toonden \textcite{strubell2019energy} aan dat het trainen van één enkel groot deep learning-model (zoals een Transformer) met \textit{Neural Architecture Search} (NAS) een $CO_2$-uitstoot kan genereren van meer dan 284.000 kg. Dit is ruwweg vergelijkbaar met de totale levensduuremissie van vijf gemiddelde auto's, inclusief de fabricage. Hoewel hun onderzoek zich primair richtte op Natural Language Processing (NLP), zijn deze bevindingen naadloos overdraagbaar op grootschalige computer vision-modellen (zoals Vision Transformers), die berucht zijn om hun rekenintensieve aard.Daarnaast bevestigen analyses van \textcite{patterson2021carbon} dat het energieverbruik van machine learning exponentieel toeneemt naarmate de datasetomvang stijgt, zelfs wanneer men optimalisaties in datacenter-infrastructuur in acht neemt.Voor organisaties vertaalt deze ecologische last zich direct in een economische last. De stijgende kosten voor cloud-infrastructuur, GPU-tijd en het opslaan van petabytes aan data creëren een aanzienlijke barrière voor innovatie. Bovendien vereisen supervised learning-modellen geannoteerde data. De tijdrovende en kapitaalintensieve aard van handmatige data-annotatie beperkt de schaalbaarheid voor kleinere bedrijven en academische instellingen, wat leidt tot een centralisatie van AI-capaciteiten bij een handvol grote technologiebedrijven.\section{Beperkingen van huidige optimalisatietechnieken}\label{sec:optimalisatie-beperkingen}De onderzoeksgemeenschap is zich bewust van deze schaalbaarheidsproblemen en heeft diverse methoden ontwikkeld om AI-modellen efficiënter te maken. Deze optimalisatietechnieken kunnen grofweg worden opgedeeld in model-centrische en data-centrische benaderingen. Geen van beide categorieën biedt echter een volledige oplossing voor het probleem van massale dataconsumptie tijdens de trainingsfase.\subsection{Model-centrische technieken}Model-centrische methoden richten zich op het verkleinen van de architectuur van het model zélf, veelal met het oog op implementatie op edge-apparaten (zoals smartphones of IoT-sensoren).\begin{itemize}\item \textbf{Pruning:} Hierbij worden de minst belangrijke verbindingen (gewichten) in een neuraal netwerk naar nul gereduceerd, wat resulteert in een \textit{sparse} model.\item \textbf{Quantization:} Het reduceren van de precisie van de gewichten, waardoor het geheugengebruik drastisch afneemt.\item \textbf{Knowledge Distillation:} Een groot, zwaar getraind 'teacher'-model wordt gebruikt om de voorspellingen van een kleiner, lichter 'student'-model te sturen.\end{itemize}Hoewel deze technieken inferentietijden drastisch versnellen en de ecologische impact tijdens het \textit{gebruik} van het model verminderen, bieden ze geen soelaas voor de trainingsfase. Om deze technieken toe te passen, moet immers eerst het grote basismodel volledig getraind worden op de complete dataset, wat betekent dat de initiële ecologische en economische kosten alsnog gemaakt worden.\subsection{Data-centrische technieken}In de verschuiving naar \textit{Data-centric AI} ligt de focus op de kwaliteit en verwerking van de data in plaats van de modelarchitectuur. Benaderingen zoals \textit{Contrastive Learning} en \textit{Self-Supervised Learning} (bijv. SimCLR \autocite{chen2020simple}) hebben de afhankelijkheid van menselijke labels aanzienlijk verminderd. Deze modellen leren robuuste representaties door het maximaliseren van de overeenkomst tussen verschillende augmentaties van dezelfde afbeelding.Desondanks hebben deze technieken een grote tekortkoming: hoewel ze de annotatiekost verlagen, verhogen ze vaak de rekenkost. Self-supervised modellen vereisen nog steeds enorme hoeveelheden ruwe, ongelabelde data om generalisatie te garanderen, en het contrastieve leerproces vergt uitzonderlijk lange trainingstijden.Dit onderzoek stelt daarom dat een werkelijk effectieve en duurzame oplossing een mechanisme vereist dat modellen in staat stelt om aan de bron selectiever te leren. In plaats van te leren van \textit{alle} beschikbare data, moet het model iteratief enkel die voorbeelden selecteren die daadwerkelijk nieuwe informatie toevoegen aan zijn huidige kennisstaat.\section{Onzekerheidskwantificatie als filter}\label{sec:onzekerheid}Om betrouwbare beslissingen te kunnen nemen over welke data informatief is en welke redundant, moet een neuraal netwerk zich bewust zijn van zijn eigen beperkingen. Standaard deterministische deep learning-modellen missen deze eigenschap: zij neigen ernaar extreem zelfverzekerd (\textit{overconfident}) te zijn in hun voorspellingen, zelfs wanneer ze geconfronteerd worden met beelden die ver buiten hun trainingsdistributie (\textit{Out-of-Distribution} of OOD) vallen. Softmax-outputs worden vaak foutief geïnterpreteerd als modelzekerheid, terwijl het in werkelijkheid slechts een genormaliseerde score is. Om een robuuste 'Smart Selection' van data te maken, moet het model echter weten \textit{wat} het niet weet.\textcite{kendall2017uncertainties} formaliseerden een raamwerk dat toelaat om onzekerheid in computer vision op te splitsen in twee fundamenteel verschillende componenten, die beide essentieel zijn voor dataselectie.\subsection{De wiskunde van Data- en Model-onzekerheid}De totale voorspellingsonzekerheid van een model is ruwweg de som van de data-onzekerheid en de model-onzekerheid.\begin{enumerate}\item \textbf{Data-onzekerheid:} Deze vorm van onzekerheid is inherent aan de waarnemingen of sensordata zelf. Voorbeelden in computer vision zijn zware compressie-artefacten, bewegingsonscherpte (\textit{motion blur}), of camera-ruis door slechte belichting. Het cruciale kenmerk van data-onzekerheid is dat het onherleidbaar is. Zelfs met een oneindig grote dataset zal de ruis in die specifieke afbeelding blijven bestaan. Wiskundig wordt dit als volgt uitgedrukt:\end{enumerate}\subsection{Bayesian Neural Networks (BNN's) voor Model-onzekerheid}Klassieke neurale netwerken representeren gewichten als vaste, deterministische waarden. Een Bayesian Neural Network (BNN) vervangt deze vaste punt-gewichten echter door waarschijnlijkheidsverdelingen. In plaats van te beweren dat een gewicht precies $0.45$ is, stelt een BNN dat het gewicht hoogstwaarschijnlijk rond $0.45$ ligt, met een zekere standaardafwijking \autocite{blundell2015weight}.Tijdens een predictie 'samplet' het netwerk een waarde uit elke distributie. De exacte Bayesiaanse voorspelling vereist het berekenen van de marginale waarschijnlijkheid, wat in de volgende formule wordt uitgedrukt:$$P(y | x, D) = \int P(y | x, \omega) P(\omega | D) d\omega$$\textit{Verduidelijking van de formule:} Deze formule beschrijft de theoretisch perfecte manier om een voorspelling te doen. Het integraal-teken ($\int$) staat voor een allesomvattende optelsom. De formule stelt dat het model, om de perfecte voorspelling te maken, elke mogelijke regel of parametercombinatie ($\omega$) die het had kunnen leren moet evalueren, en al die uitkomsten moet samenvoegen op basis van hoe waarschijnlijk die regels zijn ($P(\omega | D)$). Omdat een modern neuraal netwerk miljoenen parameters heeft, is deze berekening wiskundig onhaalbaar (\textit{intractable}). Het kost domweg te veel tijd en rekenkracht om oneindig veel netwerkcombinaties door te rekenen.Om dit op te lossen introduceerden \textcite{gal2016dropout} de \textit{Monte Carlo Dropout} (MC Dropout) techniek als een praktische benadering. Door standaard dropout niet uit te schakelen tijdens inferentie, maar juist meerdere malen toe te passen (stochastische \textit{forward passes}), creëert men een ensemble van sub-netwerken. De spreiding (variantie) in de voorspellingen van deze sub-netwerken vormt een betrouwbare maatstaf voor de model-onzekerheid \autocite{jospin2022hands}:$$\text{Uncertainty}_{model} \approx \frac{1}{T} \sum_{t=1}^{T} (\hat{y}_t - \bar{y})^2$$\textit{Verduidelijking van de formule:} Omdat de perfecte berekening uit de vorige formule onmogelijk is, functioneert deze formule als een slimme simulatie. Het sommatieteken ($\sum$) geeft aan dat het model het simpelweg een aantal keer ($T$) opnieuw probeert te raden. Bij elke poging ($t$) wordt een ander willekeurig deel van het 'brein' van het netwerk tijdelijk uitgezet (dropout), wat resulteert in een gok ($\hat{y}_t$). Vervolgens wordt berekend hoever elke individuele gok afwijkt van de gemiddelde gok ($\bar{y}$). Dit verschil wordt gekwadrateerd ($^2$).Als het model een beeld heel goed kent, zullen alle sub-netwerken het grotendeels met elkaar eens zijn en liggen de antwoorden dicht bij elkaar (lage variantie/onzekerheid). Behoort het beeld tot de blinde vlekken van het model, dan zullen de uitgeschakelde neuronen grote schommelingen in de voorspellingen veroorzaken (hoge variantie/onzekerheid).\subsection{Interval Neural Networks (INN's) voor Data-onzekerheid}Hoewel BNN's uitstekend presteren voor model-onzekerheid, bieden Interval Neural Networks (INN's) een krachtig, wiskundig alternatief om specifiek data-onzekerheid te isoleren. INN's verwerken inputdata en gewichten als wiskundige intervallen $[x_{min}, x_{max}]$ in plaats van discrete scalairen \autocite{ishibuchi1993architecture}.Wanneer een afbeelding ruis of compressiefouten bevat, kan de pixelwaarde worden beschouwd als een waarde die binnen een bepaalde foutmarge valt. Deze intervallen worden door de lagen van het netwerk gepropageerd via interval-rekenkunde. Voor een simpele optelling van twee intervallen geldt bijvoorbeeld:$$[\underline{a}, \overline{a}] + [\underline{b}, \overline{b}] = [\underline{a} + \underline{b}, \overline{a} + \overline{b}]$$\textit{Verduidelijking van de formule:} In plaats van te rekenen met exacte getallen, rekent interval-rekenkunde met marges (van-tot). De formule toont aan hoe eenvoudig dit werkt: bij een optelling neem je simpelweg de absoluut laagste waarden ($\underline{a}$ en $\underline{b}$) om het nieuwe minimum te berekenen, en de absoluut hoogste waarden ($\overline{a}$ en $\overline{b}$) voor het nieuwe maximum.Een INN past dit concept toe op het hele netwerk en resulteert in een classificatie-interval $[\underline{y}, \overline{y}]$. De breedte van dit output-interval ($\overline{y} - \underline{y}$) is een directe kwantificering van de data-onzekerheid. Hoe groter de fundamentele ruis in de inputafbeelding, hoe breder dit interval zich door het netwerk zal verspreiden. Wanneer de output-bandbreedte extreem breed wordt, signaleert het INN dat de inputdata te veel inherente ruis bevat om betrouwbaar te zijn.\section{Toepassing: De 'Smart Selection'-methode via Actief Leren en Coresets}\label{sec:smart-selection}De bovengenoemde theorieën rondom scaling laws en onzekerheidskwantificatie komen in deze bachelorproef samen in de ontwikkeling van de 'Smart Selection'-methode. Het doel van deze methode is niet om modellen na de training te verkleinen, maar om de datastroom naar het model toe dynamisch en intelligent te filteren.Dit concept leunt op de theorie van \textit{Active Learning} (Actief Leren) en \textit{Coreset Selection}. Binnen actief leren observeert een model een pool van ongelabelde of ongeziene data en 'vraagt' het actief om de annotatie (of opname in de trainingsset) van uitsluitend die beelden waar het onzeker over is \autocite{settles2009active}. Veelal gebruikt men hierbij acquisitiefuncties die direct gebaseerd zijn op de model-onzekerheid gegenereerd door BNN's \autocite{houlsby2011bayesian}.Daarnaast biedt coreset-theorie een meetkundig perspectief. Een coreset is een verkleinde subset van de originele dataset die de eigenschap heeft dat het trainen van een model op deze subset nagenoeg dezelfde test-loss oplevert als trainen op de gehele dataset \autocite{sener2017active, mirzasoleiman2020coresets}.Door deze domeinen te fuseren, past de 'Smart Selection'-methode in deze thesis een drievoudige filterstrategie toe:\begin{itemize}\item \textbf{Behouden (Hoge Model-onzekerheid):} Beelden waar de MC Dropout of INN-variantie aangeeft dat het model twijfelt vanwege een gebrek aan kennis. Deze beelden vormen de coreset en bevatten maximale leerwaarde.\item \textbf{Verwijderen (Lage Onzekerheid):} Beelden waarbij het model nagenoeg nergens twijfelt. Deze zijn redundant; ze bevestigen wat het netwerk al met hoge probabiliteit weet en leveren minimale gradiënt-updates op.\item \textbf{Verwijderen (Hoge Data-onzekerheid):} Beelden waarbij de fout ligt in de waarneming (bijv. totale ruis). Het netwerk mag niet geforceerd worden deze beelden te 'memoriseren', aangezien dit leidt tot overfitting en degradatie van representaties.\end{itemize}De leidende hypothese van dit onderzoek is dat het consequent toepassen van dit selectiemechanisme kan resulteren in een reductie van 30\% tot 40\% van de originele dataset. Indien correct geïmplementeerd, zou de generalisatiekracht van het model gewaarborgd moeten blijven, terwijl de benodigde trainingstijd en daarmee de ecologische voetafdruk drastisch dalen.